\scanpage{5}
$(\widehat X,I)$, où $\widehat X$ est un schéma projectif, dont tte composante irréductible est de dim $1$, et dont l'ensemble singulier est formé de pts quadratiques ordinaires -- et $I$ est un sous-schéma fini étale de $\widehat X^{\mathrm{lisse}}$, ou ce qui revient au $=$, une partie finie de $\widehat X^{\mathrm{lisse}}(k)$.

Soit $\widetilde{\widehat X}=\widehat{\widetilde X}$ le normalisé de $\widehat X$, \deleted{alors}
\[
\widetilde I=\widehat{\widetilde X}\setminus\widetilde X
=\text{Image inverse de }I\text{ par }\widetilde{\widehat X}\to\widehat X
\]
l'ens. des \deleted{ses} ``points à l'infini'' de $\widetilde X$, on a
\[
(2)\qquad\widetilde I\xrightarrow{\sim}I
\]
et nous identifions généralement $\widetilde I$ et $I$, et plus généralement \deleted{$\unsure{\widehat X^{\mathrm{lisse}}}$ \unsure{à}} les pts de $\widehat X^{\mathrm{lisse}}$ aux pts correspondants de $\widetilde{\widehat X}$. Soit $\widetilde A$ l'image inverse de $A$ dans $\widetilde X$ (ou dans $\widehat{\widetilde X}$), alors
\[
(3)\qquad\widetilde A\longrightarrow A
\]
est une application de degré $2$, d'où une involution sans pts fixe $\sigma_{\widetilde A}$ dans $\widetilde A$.

Ainsi, à la courbe standart $X$ nous avons associé les systèmes de données
\mspage{4}
